{\small\setlength{\tabcolsep}{4pt}
\begin{longtable}{@{}>{\raggedright\arraybackslash}p{\dimexpr 0.0261\linewidth-2\tabcolsep\relax}>{\raggedright\arraybackslash}p{\dimexpr 0.2941\linewidth-2\tabcolsep\relax}>{\raggedright\arraybackslash}p{\dimexpr 0.3529\linewidth-2\tabcolsep\relax}>{\raggedright\arraybackslash}p{\dimexpr 0.3268\linewidth-2\tabcolsep\relax}@{}}
\caption{What this article adds, item by item: what was established before, what is new here, and the evidence in this article.}\label{tab:novelty}\\
\toprule
\# & Established before & New in this article & Evidence here\\\midrule\endfirsthead
\toprule \# & Established before & New in this article & Evidence here\\\midrule\endhead
1 & Field capacity has intrinsic definitions from the hydraulic parameters \citep{Twarakavi2009,Assouline2014}; uptake becomes limited by the conductance of the soil around the roots \citep{Gardner1960,Carminati2020}. & A closed form for the dry-end counterpart, the \emph{reach} of a suction trigger, with one constant shared by all soils. & Held-out soils within a factor of 2.3 (LOSO 0.18 decades); a threshold on absolute conductivity rejected.\\
2 & Triggers of 300 to 600\,cm are tabulated for field crops \citep{Allen1998}. & These triggers lie beyond the reach of coarse soils: they never fire and the crop stays rainfed. & Sand and loamy sand at 300\,cm: rainfed yield; a deeper sensor does not help (E3).\\
3 & -- & A reach-scaled trigger $\htr/\hc$ that places the irrigation responses of all soils on one axis. & Between-soil scatter falls from 54\,\% to 32\,\% of the total.\\
4 & Pulsing, discharge and hydraulic properties change losses below the root zone, shown case by case \citep{Elmaloglou2007,Elmaloglou2008,Elmaloglou2009b}. & The marginal fate of an irrigation millimetre across soils, doses, seasons and methods. & 0.38 transpired, 0.25 soil evaporation, 0.19 other evaporation, 0.16 drained.\\
5 & -- & A dose scale of the drained share that does not depend on the soil, and a soil-specific plateau. & $d_D=4.1\pm0.6$\,mm, $\rho=0.07$ with $n$; plateau 0.10 to 0.45.\\
6 & Interception losses of sprinklers are known. & Overhead irrigation adds no transpiration: intercepted water replaces it. & $\eta_T=-0.08$ (all soils) and $-0.01$ (coarse loam).\\
7 & Flow can be out of equilibrium with the retention curve \citep{RossSmettem2000,Diamantopoulos2012VZJ}. & A criterion, $\Pi_\tau=\tau/t_i$, for when this changes the firing of a trigger (proposed, not tested). & Derivation in Section \ref{der:s4:8}.\\
\bottomrule
\end{longtable}}
